\documentclass[11pt]{article}
\usepackage[margin=1in]{geometry}
\usepackage{booktabs}
\usepackage{array}
\usepackage{xcolor}
\usepackage{hyperref}
\usepackage{amsmath}
\usepackage{amssymb}
\usepackage{enumitem}
\usepackage{titlesec}
\usepackage{longtable}
\usepackage{parskip}

\definecolor{linknavy}{HTML}{2a4d8f}
\hypersetup{colorlinks=true, linkcolor=linknavy, urlcolor=linknavy, citecolor=linknavy}

\titleformat{\section}{\Large\bfseries}{\thesection}{0.6em}{}
\titleformat{\subsection}{\large\bfseries}{\thesubsection}{0.6em}{}

\setlist[itemize]{topsep=2pt, itemsep=2pt, parsep=0pt}

\newcommand{\Nb}{N}
\newcommand{\nk}{n_k}

\title{Vectorizations for Multi-Scale Persistence Regression\\
\large Persistence Images, Landscapes, Silhouettes, and Statistics in \texorpdfstring{\texttt{vec\_multik}}{vec\_multik}}
\author{point-process-tda}
\date{\today}

\begin{document}
\maketitle
\vspace{-2em}
\tableofcontents

\section{Introduction to each vectorization}
\label{sec:intro}

Every vectorization below turns one persistence diagram (a finite set of $(b_i,d_i)$
birth--death pairs for one homology dimension, one DTM scale) into a fixed-size
numeric object a neural encoder can read. Three of the four ($\lambda_k$, $\phi^{(p)}$,
and, indirectly, the persistence image) are built from the same underlying object: a
\emph{tent function}
\begin{equation}
  \Lambda_i(t) \;=\; \max\bigl(0,\ \min(t-b_i,\ d_i-t)\bigr),
  \label{eq:tent}
\end{equation}
a piecewise-linear bump that is zero outside $[b_i,d_i]$, rises linearly from $b_i$,
peaks at the pair's midpoint $\tfrac{b_i+d_i}{2}$ with height $\tfrac{d_i-b_i}{2} =
\tfrac{p_i}{2}$ (half the persistence $p_i \equiv d_i - b_i$), and falls linearly back
to zero at $d_i$.

\subsection{Persistence image (PI)}
The established baseline in this codebase (\texttt{pi\_multik}), following Adams et
al.\ (2017). Each point becomes an isotropic 2-D Gaussian centered at
$(b_i, p_i)$ in \emph{birth--persistence} coordinates, linearly weighted by $p_i$ (zero
weight exactly on the diagonal), and the raster is the exact closed-form integral of
that Gaussian mixture over each pixel's box, not a point sample at the pixel center.
Captures \emph{where in (birth, persistence) space} mass concentrates, smoothed by a
fixed-bandwidth kernel.

\subsection{Persistence landscape}
Built directly from Equation~\ref{eq:tent} (Bubenik, 2015). The $k$-th landscape
function is the $k$-th largest tent value at each grid point $t$:
\begin{equation}
  \lambda_k(t) \;=\; k\text{-th largest of } \{\Lambda_i(t)\}_{i=1}^{M},
  \qquad \lambda_k(t) := 0 \text{ if fewer than } k \text{ tents are nonzero at } t.
  \label{eq:landscape}
\end{equation}
$\lambda_1(t)$ is the pointwise upper envelope of every tent; $\lambda_2, \lambda_3,
\ldots$ peel off successively smaller order statistics, so stacking $K$ layers exposes
\emph{how many} features are simultaneously alive at each $t$ -- a multiplicity signal
a single persistence-image raster does not expose directly.

\subsection{Persistence silhouette}
Also built from Equation~\ref{eq:tent} (Chazal et al., 2014), but instead of order
statistics it is a persistence-weighted \emph{average} tent height:
\begin{equation}
  \phi^{(p)}(t) \;=\; \frac{\sum_i |d_i-b_i|^p\, \Lambda_i(t)}{\sum_i |d_i-b_i|^p}.
  \label{eq:silhouette}
\end{equation}
The exponent $p$ is an \emph{experimental axis}, not a tuned hyperparameter: $p=0$
gives an unweighted mean over every tent (a fleeting, likely-noise feature counts
exactly as much as a long-lived one); larger $p$ increasingly emphasizes long-lived
features; $p\to\infty$ degenerates to a one-hot weighting on the single
longest-persistence pair, i.e.\ $\phi^{(\infty)}(t) = \Lambda_{i^\ast}(t)$ for
$i^\ast = \arg\max_i (d_i-b_i)$. Every $p\in\{0,1,2,3\}$ is run as a separate arm and
reported; none is selected as a ``winner.'' This build also keeps the
\emph{unnormalized} numerator of Equation~\ref{eq:silhouette} as a second channel
alongside $\phi^{(p)}$ -- the denominator discards total feature count/weight, which
plausibly matters for recovering intensity-like parameters.

\subsection{Persistence statistics (floor control)}
No spatial or rank structure at all: 18 scalars per (dimension, scale) -- point count,
total persistence $\sum_i(d_i-b_i)$, persistent entropy (Shannon entropy of the
normalized lifetime distribution), and the $\{10,25,50,75,90\}$-th percentiles of
birth, death, and lifetime separately. Exists to answer whether the raster/curve
\emph{shape} the other three vectorizations preserve is actually earning its keep over
a handful of order statistics.

\section{High-level pipeline used for each}
\label{sec:pipeline}

All four sit inside the same outer architecture, \texttt{vec\_multik}
(\texttt{experiments/pi\_multik/vectorized\_multik.py}), which generalizes the
original \texttt{pi\_multik} late-fusion design across a \texttt{vectorization} axis
and an \texttt{encoder\_path} axis:

\begin{enumerate}
  \item \textbf{Sample} a point cloud from a process (Thomas, nested Thomas) over the
    unit square.
  \item \textbf{Filter} at one or more DTM neighborhood scales $k \in \{5,10,15\}$
    (or all three fused at once), producing one persistence diagram ($H_0$ and $H_1$
    birth--death pairs) per scale.
  \item \textbf{Vectorize} each (dimension, scale) diagram independently into a
    fixed-size tensor -- this step is the one thing that differs across the four
    vectorizations (Table~\ref{tab:pervec}).
  \item \textbf{Encode} each scale's tensor with one encoder into a fixed-width
    embedding. \texttt{EncoderBank} handles the scale axis identically regardless of
    vectorization: with the default \texttt{encoder\_mode=shared}, one encoder
    instance is applied to every scale by folding the scale axis into the batch
    dimension (one call per forward pass, weight-tied across scales), rather than one
    encoder per scale.
  \item \textbf{Fuse} the resulting $(\nk, E)$ embedding sequence into one vector --
    flat concatenation by default (\texttt{fusion\_mode=concat}), giving a
    $\nk\!\cdot\!E$-dimensional vector; an optional \texttt{Conv1d}-over-scales fusion
    (\texttt{ConvFusion}) exists but is not used in the results referenced here.
  \item \textbf{Append} scalar side-features -- $\log N(x)$ (point count) always,
    optionally per-(dimension, scale) persistence entropy.
  \item \textbf{Regress} the process parameters with a shared MLP head
    (\texttt{ParameterEstimator}), trained with MSE loss in log-z-score-normalized
    target space.
\end{enumerate}

Only two things vary by vectorization: what shape the per-(dimension, scale) tensor
has, and which encoder is used to read it.

\begin{table}[h]
\centering
\begin{tabular}{@{}lll@{}}
\toprule
Vectorization & Tensor per (dim, scale) & Native encoder \\
\midrule
Persistence image  & $(64,64)$ raster & \texttt{CoordConvPIEncoder} (2-D CNN, unchanged) \\
Persistence landscape & $(K,G)$ raster & same \texttt{CoordConvPIEncoder}, reused \\
Persistence silhouette & $(G,)$ curve $\times\,2$\footnotemark & \texttt{SilhouetteConv1DEncoder} (new, 1-D) \\
Persistence statistics & $(18,)$ vector & none -- MLP-only \\
\bottomrule
\end{tabular}
\caption{What differs across vectorizations. Every row also supports the shared
\texttt{FlattenMLPEncoder} path (\texttt{encoder\_path=mlp}); statistics supports
\emph{only} that path, since there is no spatial structure for a CNN to exploit.}
\label{tab:pervec}
\end{table}
\footnotetext{Normalized $\phi^{(p)}$ and its unnormalized numerator, see Section~\ref{sec:intro}.}

Both CNN-style encoders append coordinate-ramp channels before their first
convolution -- two channels ($x,y\in[-1,1]$ meshgrid ramps) for the 2-D case, one
channel (a single ramp along the grid axis) for the 1-D case -- so the conv stack has
explicit access to \emph{where} along the raster/curve it is, not just \emph{what
shape} it is looking at (plain convolutions are translation-invariant and otherwise
blind to absolute position). The \texttt{FlattenMLPEncoder} path
(\texttt{flatten}~$\to$~\texttt{Linear(512)}~$\to$~\texttt{ReLU}~$\to$~\texttt{Dropout(0.1)}~$\to$~\texttt{Linear(128)}~$\to$~\texttt{ReLU})
has no such augmentation and no notion of raster shape at all -- it is deliberately
architecture-agnostic, run identically over every vectorization (including
persistence images) as the encoder-confound-free comparison path.

One combination is not reimplemented: \texttt{vectorization=persistence\_image},
\texttt{encoder\_path=native} delegates \texttt{vec\_multik}'s entire \texttt{run()}
call straight to the original, unmodified \texttt{PIMultiKExperiment}, so it is
byte-identical to plain \texttt{pi\_multik} for the same seed -- not merely
numerically close.

\section{Calibration procedure for each}
\label{sec:calibration}

Every calibrated quantity below -- box/grid range, landscape depth $K$, or
z-score statistics -- is fit \emph{once per seed}, on that seed's TRAIN split only
(the same 70/15/15 split used everywhere in this pipeline), then frozen and applied
unchanged to that seed's validation, test, and adversarial diagrams. Two different
seeds therefore legitimately end up with different calibrated ranges; this is a
regression-tested invariant (see \texttt{tests/test\_pipeline\_e2e.py} and
\texttt{tests/test\_landscape\_silhouette.py}), not an incidental property.

\subsection{Persistence image}
A 2-D box, $\text{birth\_range}=[\,b_{\mathrm{lo}}, b_{\mathrm{hi}}\,]$,
$\text{pers\_range}=[\,0, p_{\mathrm{hi}}\,]$, fit from per-diagram extrema (minimum
birth, maximum birth, maximum persistence) pooled over training diagrams at coverage
$q=0.99$ with pad multiplier $1.05$:
\[
  b_{\mathrm{lo}} = Q_{1-q}\bigl(\{\min_i b_i\}\bigr), \qquad
  b_{\mathrm{hi}} = 1.05\cdot Q_{q}\bigl(\{\max_i b_i\}\bigr), \qquad
  p_{\mathrm{hi}} = 1.05\cdot Q_{q}\bigl(\{\max_i p_i\}\bigr),
\]
the maximum/minimum taken per diagram, the quantile $Q$ taken across diagrams. Raises
if $b_{\mathrm{hi}} \le b_{\mathrm{lo}}$ (a degenerate birth axis -- e.g.\ plain Rips
$H_0$, where every feature is born at filtration value $0$). Resolution is fixed at
$64\times64$; \texttt{sigma\_pixels} sets the Gaussian bandwidth in \emph{pixel} units
(converted to absolute birth/persistence-unit $\sigma$ by \texttt{sigma\_pixels}
$\times$ axis span $/$ resolution), so the visual smoothing is comparable even though
the two calibrated axes end up with very different physical spans.

\subsection{Persistence landscape}
A 1-D range $[t_{\min}, T]$, calibrated per (homology dimension, scale) the same way
as the birth axis above, but with no birth-axis degeneracy concern (there is no birth
axis) and no clipping (content outside $[t_{\min},T]$ is simply never evaluated):
\[
  t_{\min} = Q_{1-q}\bigl(\{\min_i b_i\}\bigr), \qquad
  T = \texttt{pad\_factor}\cdot Q_{q}\bigl(\{\max_i d_i\}\bigr), \qquad
  \texttt{grid} = \mathrm{linspace}(t_{\min}, T, G).
\]
$q=0.99$ and \texttt{pad\_factor}$=1.05$ are inherited unchanged from the
persistence-image calibration above. $H_0$ births are nonzero under DTM, so
$t_{\min}$ is always fit, never assumed to be $0$.

$K$ (landscape depth) is chosen by a coverage rule, \emph{independently per
homology dimension}: the smallest $K$ such that layer $K{+}1$'s sup-norm is
$\le 1\%$ of layer $1$'s sup-norm for at least $q$ of the training diagrams, capped at
$16$. Because $H_0$ (typically many simultaneously-alive features under DTM) and $H_1$
(typically few) routinely pick different $K$ under this rule, the two per-dimension
answers are \emph{reconciled to a single shared} $K = \max(K_{H_0}, K_{H_1})$ before
the tensor is built -- both dimensions' own coverage-rule answer is still logged, but
only the shared value is used, since $H_0$ and $H_1$ become two \emph{channels} of one
rectangular $(C,K,G)$ raster fed to a single CNN, which requires identical spatial
size on every channel. $G$, and an explicit $K$ override when one is supplied (e.g.\
the budget-matched arm's $K=8$), are config choices, not fit from data.

\subsection{Persistence silhouette}
Identical $[t_{\min},T]$ grid calibration as the landscape's, per (dimension, scale).
No $K$ axis (a silhouette is one curve, not a stack of layers). $p$ is fixed for the
whole run, chosen by the experimenter as one of the arms $\{0,1,2,3\}$ -- never
data-derived, unlike $K$.

\subsection{Persistence statistics}
No grid, box, or depth to calibrate: the 18 statistics are a pure per-diagram
function, calibration-free in the same sense as persistence entropy. The one
train-fit/apply-frozen step is a per-column z-score, applied before the MLP encoder
(mean/std of each of the 18 statistics, fit on that seed's train rows, frozen for
val/test/adversarial) -- necessary because a raw count, a filtration-value quantile,
and an entropy live on very different numeric scales, and there is no
BatchNorm-bearing conv stack downstream to absorb that automatically, unlike the
raster/curve vectorizations above (whose raw magnitude is left deliberately
unnormalized -- see Section~\ref{sec:shapes} -- since e.g.\ a larger DTM scale $k$
tends to produce larger raw persistences, and that scale information is itself
useful signal for the shared-weight encoder to exploit).

\section{Detailed pipeline: shapes at each stage}
\label{sec:shapes}

Notation: $\Nb$ = batch size, $\nk$ = number of DTM scales fused ($3$ for the combined
$k\in\{5,10,15\}$ runs, $1$ for a single-scale run), $C=2$ = number of homology
dimensions ($H_0, H_1$), $G$ = grid resolution, $K$ = landscape depth (reconciled
across dimensions, Section~\ref{sec:calibration}), $E$ = embedding dimension ($64$ for
every native-encoder config referenced here; the \texttt{FlattenMLPEncoder} path fixes
$E=128$ regardless of config), $n_{\mathrm{extra}}=1$ (just $\log N(x)$ in every config
referenced here), $n_{\mathrm{targets}} \in \{3,5\}$ (Thomas / nested Thomas). Concrete
numbers below use the actual ``native'' \texttt{vec\_multik} configs
(\texttt{configs/runs/\{thomas,nested\_thomas\}/vec\_multik\_*\_native.yaml}):
$G=128$, resolution $64\times64$ for images, $E=64$.

Every vectorization shares the same last four rows (\texttt{EncoderBank}'s batch-fold,
fusion, extra-feature concatenation, and head) -- only the first few rows, building
the per-scale tensor and reading it with the native encoder, differ.

\subsection{Persistence image}
\begin{center}
\begin{tabular}{@{}p{6.6cm}p{8.6cm}@{}}
\toprule
Stage & Shape \\
\midrule
Diagram, one (dim, scale) & $(M_i,\,2)$ birth--death pairs \\
Rasterized image, one (dim, scale) & $(64,\,64)$ \\
Stacked over $C$ dims, one scale & $(C,\,64,\,64) = (2,64,64)$ \\
Stacked over $\nk$ scales & $(\Nb,\ \nk,\ C,\ 64,\ 64)$ \\
\texttt{EncoderBank} fold (scale $\to$ batch) & $(\Nb\!\cdot\!\nk,\ C,\ 64,\ 64)$ \\
$+2$ coordinate channels & $(\Nb\!\cdot\!\nk,\ C{+}2,\ 64,\ 64) = (\cdot,4,64,64)$ \\
\texttt{Conv2d} $\times3$ (channels $4\to32\to64\to128$) + pools & $(\Nb\!\cdot\!\nk,\ 128,\ h,\ w)$, $h,w$ shrunk by two $2\!\times\!2$ pools \\
\texttt{AdaptiveAvgPool2d(1)} + flatten + \texttt{Linear(128,\,64)} & $(\Nb\!\cdot\!\nk,\ 64)$ \\
Reshape (batch $\to$ scale) & $(\Nb,\ \nk,\ 64)$ \\
Flat-concat fusion & $(\Nb,\ \nk\!\cdot\!64)$ -- $193$-d for $\nk{=}3$ \\
$+$ extra features & $(\Nb,\ \nk\!\cdot\!64 + 1)$ \\
Head: \texttt{Linear}$\to$\texttt{64}$\to$\texttt{32}$\to n_{\mathrm{targets}}$ & $(\Nb,\ n_{\mathrm{targets}})$ \\
\bottomrule
\end{tabular}
\end{center}

\subsection{Persistence landscape}
\begin{center}
\begin{tabular}{@{}p{6.6cm}p{8.6cm}@{}}
\toprule
Stage & Shape \\
\midrule
Diagram, one (dim, scale) & $(M_i,\,2)$ birth--death pairs \\
Tent matrix (Eq.~\ref{eq:tent}) & $(M_i,\,G)$ \\
Landscape, one (dim, scale) & $(K,\,G)$ -- $K$ layers, $K$ shared across $C$ (Section~\ref{sec:calibration}) \\
Stacked over $C$ dims, one scale & $(C,\,K,\,G) = (2,K,128)$ \\
Stacked over $\nk$ scales & $(\Nb,\ \nk,\ C,\ K,\ G)$ \\
\texttt{EncoderBank} fold & $(\Nb\!\cdot\!\nk,\ C,\ K,\ G)$ \\
$+2$ coordinate channels & $(\Nb\!\cdot\!\nk,\ C{+}2,\ K,\ G) = (\cdot,4,K,128)$ \\
\texttt{Conv2d} $\times3$ + pools (\emph{same architecture as PI, no longer assumed square}) & $(\Nb\!\cdot\!\nk,\ 128,\ k',\ g')$ \\
\texttt{AdaptiveAvgPool2d(1)} + flatten + \texttt{Linear(128,\,64)} & $(\Nb\!\cdot\!\nk,\ 64)$ \\
Reshape, fusion, extra, head & identical to the PI rows above \\
\bottomrule
\end{tabular}
\end{center}
The landscape raster is \emph{not} square ($K\le16 \ne G=128$ in every config used
here), which required the coordinate-channel builder to construct independent
height/width ramps from the raster's own $H$ and $W$ rather than one ramp reused for
both axes -- the latter silently worked for persistence images (always
resolution$\times$resolution) but crashes the channel concatenation the moment $K\neq
G$. Two pooling stages of stride $2$ are why $K\ge4$ is asserted at model-build time.

\subsection{Persistence silhouette}
\begin{center}
\begin{tabular}{@{}p{6.6cm}p{8.6cm}@{}}
\toprule
Stage & Shape \\
\midrule
Diagram, one (dim, scale) & $(M_i,\,2)$ birth--death pairs \\
Tent matrix (Eq.~\ref{eq:tent}) & $(M_i,\,G)$ \\
Silhouette, one (dim, scale) & $\phi^{(p)}$: $(G,)$, plus unnormalized numerator: $(G,)$ \\
Stacked over $C$ dims $\times\,2$ variants, one scale & $(2C,\,G) = (4,128)$ \\
Stacked over $\nk$ scales & $(\Nb,\ \nk,\ 2C,\ G)$ \\
\texttt{EncoderBank} fold & $(\Nb\!\cdot\!\nk,\ 2C,\ G)$ \\
$+1$ coordinate-ramp channel & $(\Nb\!\cdot\!\nk,\ 2C{+}1,\ G) = (\cdot,5,128)$ \\
\texttt{Conv1d} $\times3$ (channels $5\to32\to64\to128$) + 2 max-pools & $(\Nb\!\cdot\!\nk,\ 128,\ g')$ \\
\texttt{AdaptiveAvgPool1d(1)} + flatten + \texttt{Linear(128,\,64)} & $(\Nb\!\cdot\!\nk,\ 64)$ \\
Reshape, fusion, extra, head & identical to the PI rows above \\
\bottomrule
\end{tabular}
\end{center}

\subsection{Persistence statistics}
\begin{center}
\begin{tabular}{@{}p{6.6cm}p{8.6cm}@{}}
\toprule
Stage & Shape \\
\midrule
Diagram, one (dim, scale) & $(M_i,\,2)$ birth--death pairs \\
Statistics, one (dim, scale) & $(18,)$ -- count, total persistence, entropy, 5 quantiles $\times$ 3 (birth/death/lifetime) \\
Stacked over $C$ dims, one scale & $(C,\,18) = (2,18)$ \\
Stacked over $\nk$ scales & $(\Nb,\ \nk,\ C,\ 18)$ \\
Per-column z-score (train-fit/apply-frozen) & $(\Nb,\ \nk,\ C,\ 18)$, unchanged shape \\
\texttt{EncoderBank} fold & $(\Nb\!\cdot\!\nk,\ C,\ 18)$ \\
Flatten & $(\Nb\!\cdot\!\nk,\ 36)$ \\
\texttt{Linear(36,512)}$\to$\texttt{ReLU}$\to$\texttt{Dropout(0.1)}$\to$\texttt{Linear(512,128)}$\to$\texttt{ReLU} & $(\Nb\!\cdot\!\nk,\ 128)$ \\
Reshape & $(\Nb,\ \nk,\ 128)$ \\
Flat-concat fusion & $(\Nb,\ \nk\!\cdot\!128)$ -- $384$-d for $\nk{=}3$ \\
$+$ extra, head & $(\Nb,\ \nk\!\cdot\!128+1) \to (\Nb,\ n_{\mathrm{targets}})$ \\
\bottomrule
\end{tabular}
\end{center}
Persistence statistics is the only vectorization with \emph{no} native-encoder row --
it is always read by the same \texttt{FlattenMLPEncoder} every other vectorization can
optionally use instead of its own native CNN, which is why $E=128$ here rather than
the $64$ every native path in this document uses.

\end{document}
